% GENERATED FILE. Edit canonical lessons, then run npm run generate:books.
% canonical-source-hash: fnv1a64:601a03534517569b
% canonical-lessons: GU-C260-pattu, GU-C260-shatranj, GU-C260-ramakdu, GU-C260-dharvu, GU-C260-pasand-karvu, GU-R260-review-255-260

\chapter{Cards, chess and toys}
\label{ch:gu-cards-chess-and-toys}
\hlchaptermodality{260}

\begin{chapteropening}
\textbf{By the end of this chapter:} \emph{I can say which games I chose and what I guessed.}

Complete the last lesson of chapter 260: I can say which games I chose and what I guessed.
\end{chapteropening}

\section[\texorpdfstring{\gu{પત્તું} (pattũ)}{પત્તું (pattũ)}]{\gu{પત્તું} (pattũ) — playing card}

\label{lesson:GU-C260-pattu}



\begin{quote}
Before the new one: say the Gujarati for exercise, then the Gujarati for wrestling.
\end{quote}

\subsection*{You'll want to know: \gu{પત્તું}}
\textbf{\gu{પત્તું}} — \emph{pattũ} — \textquotedblleft{}playing card\textquotedblright{}.

\begin{quote}
\textbf{\gu{અમે} \gu{રાત્રે} \gu{પત્તાં} \gu{રમ્યા}.} — \emph{ame rātre pattā̃ ramyā} — \textquotedblleft{}We played cards at night.\textquotedblright{}
\end{quote}

The plural is \textbf{\gu{પત્તાં}}, \emph{pattā̃}, \textquotedblleft{}cards\textquotedblright{}.

\begin{cousinweb}
\textbf{\gu{પત્તું}} (\emph{pattũ}), a leaf, then a playing card, is a tadbhava. It comes from Sanskrit \emph{páttra}, \textquotedblleft{}wing-feather; leaf; leaf of a book\textquotedblright{}; CDIAL lists the Gujarati \emph{pātũ}, \textquotedblleft{}leaf\textquotedblright{}, but not the form \emph{pattũ} (CDIAL 7733).
\end{cousinweb}

\begin{grammarlens}[title={What you've built}]
One more word for everyday life at A2.
\end{grammarlens}

\subsection*{Your turn}
\begin{itemize}
\raggedright
  \item \emph{Say it:} \emph{pattũ}
  \item \emph{Say it:} \emph{pattũ}, once more
  \item \emph{Say it:} \emph{ame rātre pattā̃ ramyā} — We played cards at night.
  \item \emph{Recall:} say the Gujarati for exercise, then the Gujarati for wrestling, then say \emph{pattũ} again
\end{itemize}

\begin{tcolorbox}[breakable,colback=teal!4,colframe=teal!35!black,title={Before you move on}]
What does \gu{પત્તું} mean? (\textquotedblleft{}Playing card\textquotedblright{}.) Say it once more.
\end{tcolorbox}
\section[\texorpdfstring{\gu{શતરંજ} (śatranj)}{શતરંજ (śatranj)}]{\gu{શતરંજ} (śatranj) — chess}

\label{lesson:GU-C260-shatranj}



\begin{quote}
Before the new one: say the Gujarati for wrestling, then the Gujarati for playing card.
\end{quote}

\subsection*{You'll want to know: \gu{શતરંજ}}
\textbf{\gu{શતરંજ}} — \emph{śatranj} — \textquotedblleft{}chess\textquotedblright{}.

\begin{quote}
\textbf{\gu{દાદાએ} \gu{મને} \gu{શતરંજ} \gu{શીખવી}.} — \emph{dādāe mane śatranj śīkhavī} — \textquotedblleft{}Grandfather taught me chess.\textquotedblright{}
\end{quote}

A feminine noun.

\begin{cousinweb}
\textbf{\gu{શતરંજ}} (\emph{śatranj}) is a loan from Persian \emph{shatranj}, \textquotedblleft{}chess\textquotedblright{} (EtymDB); Lane records the Arabic \emph{shiṭranj} as a Persian word taken into Arabic (Lane). The Persian word comes from Sanskrit \emph{cáturaṅga}, \textquotedblleft{}having four limbs\textquotedblright{}, the four parts of an army, which CDIAL also glosses \textquotedblleft{}a kind of chess\textquotedblright{} (CDIAL 4596).
\end{cousinweb}

\begin{grammarlens}[title={What you've built}]
One more word for everyday life at A2.
\end{grammarlens}

\subsection*{Your turn}
\begin{itemize}
\raggedright
  \item \emph{Say it:} \emph{śatranj}
  \item \emph{Say it:} \emph{śatranj}, once more
  \item \emph{Say it:} \emph{dādāe mane śatranj śīkhavī} — Grandfather taught me chess.
  \item \emph{Recall:} say the Gujarati for wrestling, then the Gujarati for playing card, then say \emph{śatranj} again
\end{itemize}

\begin{tcolorbox}[breakable,colback=teal!4,colframe=teal!35!black,title={Before you move on}]
What does \gu{શતરંજ} mean? (\textquotedblleft{}Chess\textquotedblright{}.) Say it once more.
\end{tcolorbox}
\section[\texorpdfstring{\gu{રમકડું} (ramakḍũ)}{રમકડું (ramakḍũ)}]{\gu{રમકડું} (ramakḍũ) — toy}

\label{lesson:GU-C260-ramakdu}



\begin{quote}
Before the new one: say the Gujarati for playing card, then the Gujarati for chess.
\end{quote}

\subsection*{You'll want to know: \gu{રમકડું}}
\textbf{\gu{રમકડું}} — \emph{ramakḍũ} — \textquotedblleft{}toy\textquotedblright{}.

\begin{quote}
\textbf{\gu{મેં} \gu{બાળક} \gu{માટે} \gu{રમકડું} \gu{ખરીદ્યું}.} — \emph{mẽ bāḷak māṭe ramakḍũ kharīdyũ} — \textquotedblleft{}I bought a toy for the child.\textquotedblright{}
\end{quote}

\begin{cousinweb}
\textbf{\gu{રમકડું}} (\emph{ramakḍũ}), a thing to play with, is built on the taught \emph{ramvũ}, \textquotedblleft{}to play\textquotedblright{}, with \emph{-kḍũ}. CDIAL derives \emph{ramvũ} from Sanskrit \emph{rámatē}, \textquotedblleft{}takes pleasure in\textquotedblright{} (CDIAL 10628), but lists no \emph{ramakḍũ}.
\end{cousinweb}

\begin{grammarlens}[title={What you've built}]
One more word for everyday life at A2.
\end{grammarlens}

\subsection*{Your turn}
\begin{itemize}
\raggedright
  \item \emph{Say it:} \emph{ramakḍũ}
  \item \emph{Say it:} \emph{ramakḍũ}, once more
  \item \emph{Say it:} \emph{mẽ bāḷak māṭe ramakḍũ kharīdyũ} — I bought a toy for the child.
  \item \emph{Recall:} say the Gujarati for playing card, then the Gujarati for chess, then say \emph{ramakḍũ} again
\end{itemize}

\begin{tcolorbox}[breakable,colback=teal!4,colframe=teal!35!black,title={Before you move on}]
What does \gu{રમકડું} mean? (\textquotedblleft{}Toy\textquotedblright{}.) Say it once more.
\end{tcolorbox}
\section[\texorpdfstring{\gu{ધારવું} (dhārvũ)}{ધારવું (dhārvũ)}]{\gu{ધારવું} (dhārvũ) — to guess, to suppose}

\label{lesson:GU-C260-dharvu}



\begin{quote}
Before the new one: say the Gujarati for chess, then the Gujarati for toy.
\end{quote}

\subsection*{You'll want to know: \gu{ધારવું}}
\textbf{\gu{ધારવું}} — \emph{dhārvũ} — \textquotedblleft{}to guess, to suppose\textquotedblright{}.

\begin{quote}
\textbf{\gu{મેં} \gu{સાચો} \gu{જવાબ} \gu{ધાર્યો}.} — \emph{mẽ sāco javāb dhāryo} — \textquotedblleft{}I guessed the right answer.\textquotedblright{}
\end{quote}

\begin{cousinweb}
\textbf{\gu{ધારવું}} (\emph{dhārvũ}), to hold a thought, is a tadbhava. It comes from Sanskrit \emph{dhāráyati}, \textquotedblleft{}holds, keeps\textquotedblright{}; CDIAL glosses the Gujarati verb \textquotedblleft{}to consider, guess\textquotedblright{} (CDIAL 6791).
\end{cousinweb}

\begin{grammarlens}[title={What you've built}]
One more word for everyday life at A2.
\end{grammarlens}

\subsection*{Your turn}
\begin{itemize}
\raggedright
  \item \emph{Say it:} \emph{dhārvũ}
  \item \emph{Say it:} \emph{dhārvũ}, once more
  \item \emph{Say it:} \emph{mẽ sāco javāb dhāryo} — I guessed the right answer.
  \item \emph{Recall:} say the Gujarati for chess, then the Gujarati for toy, then say \emph{dhārvũ} again
\end{itemize}

\begin{tcolorbox}[breakable,colback=teal!4,colframe=teal!35!black,title={Before you move on}]
What does \gu{ધારવું} mean? (\textquotedblleft{}To guess, to suppose\textquotedblright{}.) Say it once more.
\end{tcolorbox}
\section[\texorpdfstring{\gu{પસંદ} \gu{કરવું} (pasand karvũ)}{પસંદ કરવું (pasand karvũ)}]{\gu{પસંદ} \gu{કરવું} (pasand karvũ) — to choose, to pick}

\label{lesson:GU-C260-pasand-karvu}



\begin{quote}
Before the new one: say the Gujarati for toy, then the Gujarati for to guess.
\end{quote}

\subsection*{You'll want to know: \gu{પસંદ} \gu{કરવું}}
\textbf{\gu{પસંદ} \gu{કરવું}} — \emph{pasand karvũ} — \textquotedblleft{}to choose, to pick\textquotedblright{}.

\begin{quote}
\textbf{\gu{તેણે} \gu{લાલ} \gu{રમકડું} \gu{પસંદ} \gu{કર્યું}.} — \emph{teṇe lāl ramakḍũ pasand karyũ} — \textquotedblleft{}She chose the red toy.\textquotedblright{}
\end{quote}

\textbf{\gu{પસંદ}}, \emph{pasand}, alone means liked or chosen.

\begin{cousinweb}
\textbf{\gu{પસંદ} \gu{કરવું}} (\emph{pasand karvũ}) is built from \emph{pasand}, \textquotedblleft{}liked, approved\textquotedblright{}, a loan usually traced to Persian \emph{pasand}, and the verb \emph{karvũ}, \textquotedblleft{}to do\textquotedblright{}.
\end{cousinweb}

\begin{grammarlens}[title={What you've built}]
One more word for everyday life at A2.
\end{grammarlens}

\subsection*{Your turn}
\begin{itemize}
\raggedright
  \item \emph{Say it:} \emph{pasand karvũ}
  \item \emph{Say it:} \emph{pasand karvũ}, once more
  \item \emph{Say it:} \emph{teṇe lāl ramakḍũ pasand karyũ} — She chose the red toy.
  \item \emph{Recall:} say the Gujarati for toy, then the Gujarati for to guess, then say \emph{pasand karvũ} again
\end{itemize}

\begin{tcolorbox}[breakable,colback=teal!4,colframe=teal!35!black,title={Before you move on}]
What does \gu{પસંદ} \gu{કરવું} mean? (\textquotedblleft{}To choose, to pick\textquotedblright{}.) Say it once more.
\end{tcolorbox}
\section[(dialogue)]{Words from chapters 255–260 — a review}

\label{lesson:GU-R260-review-255-260}



\begin{quote}
Say the Gujarati for story and poem.
\end{quote}

\subsection*{Your turn}
Say these aloud:

\begin{itemize}
\raggedright
  \item story, poem, novel, drama, film
  \item picture, to paint, paintbrush, photo, to knit
  \item game, player, competition, to take part, to join
  \item cricket, ball, bat, to hit, to bounce
  \item match, to defeat, to run, exercise, wrestling
  \item playing card, chess, toy, to guess, to choose
\end{itemize}

\begin{tcolorbox}[breakable,colback=teal!4,colframe=teal!35!black,title={Before you move on}]
Story? (\textbf{\gu{વાર્તા}}.) Poem? (\textbf{\gu{કવિતા}}.) Novel? (\textbf{\gu{નવલકથા}}.) Drama? (\textbf{\gu{નાટક}}.) Film? (\textbf{\gu{ફિલ્મ}}.) Picture? (\textbf{\gu{ચિત્ર}}.) To paint? (\textbf{\gu{ચીતરવું}}.) Paintbrush? (\textbf{\gu{પીંછી}}.) Photo? (\textbf{\gu{ફોટો}}.) To knit? (\textbf{\gu{ગૂંથવું}}.) Game? (\textbf{\gu{રમત}}.) Player? (\textbf{\gu{ખેલાડી}}.) Competition? (\textbf{\gu{હરીફાઈ}}.) To take part? (\textbf{\gu{ભાગ} \gu{લેવો}}.) To join? (\textbf{\gu{જોડાવું}}.) Cricket? (\textbf{\gu{ક્રિકેટ}}.) Ball? (\textbf{\gu{દડો}}.) Bat? (\textbf{\gu{બેટ}}.) To hit? (\textbf{\gu{મારવું}}.) To bounce? (\textbf{\gu{ઊછળવું}}.) Match? (\textbf{\gu{મેચ}}.) To defeat? (\textbf{\gu{હરાવવું}}.) To run? (\textbf{\gu{દોડવું}}.) Exercise? (\textbf{\gu{કસરત}}.) Wrestling? (\textbf{\gu{કુસ્તી}}.) Playing card? (\textbf{\gu{પત્તું}}.) Chess? (\textbf{\gu{શતરંજ}}.) Toy? (\textbf{\gu{રમકડું}}.) To guess? (\textbf{\gu{ધારવું}}.) To choose? (\textbf{\gu{પસંદ} \gu{કરવું}}.)
\end{tcolorbox}
